\chapter{}
《魔法少女小圆》（以下简称为《魔圆》，以区别于主角鹿目圆\cite{ref1}, 后者在本文中简称“小圆”）是多重含义的。最显而易见的，它是由动画公司Shaft制作的魔法少女题材的动漫系列，但除此之外，它也是对该类型的批判和致敬；是一个迷失在黑暗中的作家自我救赎的尝试；是对浮士德的佛教式重述。它是一项对比研究，既是对同意、特权和性别角色等问题的女性主义考察，同时也是对未成年女性形象的性别歧视的客体化。它包含了牺牲与自私、绝望与希望、爱与死亡、腐朽与生命。当然，这些看似对立的概念实则是统一的。

但最重要的是，它是关于人的故事，这些人做出选择（通常是错误的），然后受苦、奋斗、成功、失败。这是关于角色们彼此相爱、相恨、相漠不关心，有时甚至皆有的故事。

在这本书中，我将讨论宏大的主题、哲学和符号，因此在开始前，我想提醒大家：这是一个关于人的故事，因为所有故事本质上都是关于人类的\cite{ref2} 。 
\paragraph{为什么研究《魔圆》？}~{}

在短短每集只有22分钟的十二集，以及一部电影中（虽然有三部电影上映，但前两部是TV动画的总集篇，只有少量细微的改动），《魔圆》呈现了极其丰富的故事，刻画了饱满的角色、复杂的主题以及极为悲伤、恐惧、愤怒和喜悦的瞬间。对它进行分析极具成就感，每次仔细研究后都能揭示出更多层次。尤其是其作为一部后现代主义作品，非常适合用后现代主义的视角进行分析。
\paragraph{“后-”什么？}~{}

“后现代主义”这个术语的定义相当复杂，但至少在谈到它时，我一定会尝试去解释清楚。在哲学上，顾名思义，后现代主义是现代主义的进一步延伸。

简而言之，现代主义的核心认识是，符号本质上是任意的。也就是说，符号与其所代表的事物之间没有固有的联系，只有观看符号的人脑海中的关联。例如，一个红色六边形意味着“停”，并不是因为红色、六边形与停止之间有某种逻辑关联，而是因为有人决定红色六边形应该代表“停”，并让其他人接受了这一点。这就是社会建构主义的核心概念，即我们理解世界的概念和符号是由社会环境和关系构建而成的。因此，这些符号和社会一样是流动的，没有固定的客观意义\cite{ref3}。

既然所有艺术（和语言）在后现代主义思想中都是一系列符号，那么任何作品的意义在某种程度上都是任意的，是社会建构的产物，而不是逻辑必然关系的结果。传统上，艺术通过使用约定俗成的文化符号来解决这一问题，如语言或具象图像（例如用苹果的图像代表一个苹果）。这些符号作为“导轨”，帮助观众从一些熟悉的符号开始接触艺术，再逐步构建出更多意义。现代主义则以拒绝传统和表现失落感、解体感为主要特征，现代艺术经常抛弃部分或全部“导轨”，激烈挑战艺术是否应该拥有非任意意义的概念。例如，詹姆斯·乔伊斯的作品以意识流的形式呈现，打破了正常的句子结构；蒙德里安的画作则用抽象的色块和形状代替了熟悉的物体。然而，现代主义常常试图重新创造它认为已经失去的秩序\cite{ref4}。例如，乔伊斯很好地将一个普通的日子描绘成一场史诗般的斗争，蒙德里安则将他的抽象色块安排成网格状的图案。

与此相反，后现代主义拒绝将社会建构的意义等同于“没有意义”。现代主义说：“符号与其所指的关系是任意的，因此一切都没有意义。”而后现代主义则说：“符号与其所指的关系是任意的，因此我们可以自由决定一切的意义。”因此，后现代主义作品不需要重新引入失去的秩序，因为那种秩序从未真正存在过。正因为如此，后现代时期对流行文化作为艺术表达的媒介的接受度也更高\cite{ref5}。现代主义拒绝传统的意义、形式和技巧，以便创造一种新的秩序来取代旧的；而后现代主义则将传统与新事物结合起来，以追求怪异、有趣的东西。\cite{ref6}

现代主义通过消除“导轨”，而后现代主义则通过强调“导轨”来创作。这主要通过去情境化和再情境化的过程实现，将熟悉的符号从其通常的情境中移除，并将其放置在其他情境中，例如通过混合多种类型的元素，创作出常见作品的集锦，或打破“第四面墙”\cite{ref7}。这种引起的迷失感通常用于制造幽默或恐怖效果，或者引发思考（当然，也可以同时既有趣又可怕，并且发人深省）。

后现代主义可以分为“积极”倾向，强调玩乐和自由，以及“消极”倾向，强调混乱和犬儒主义或绝望\cite{ref8}。《魔圆》很大程度上属于后者，并大量运用了后现代主义技巧来讲述它的故事，例如（我们将在下一章更详细讨论）让“魔女”的形象采用外来、侵略性的艺术风格，取代了其熟悉的动画风格。
\paragraph{关于本书}~{}

本书分为三个部分。第一部分是TV动画，包含十二章，每一章分别讲述TV动画《魔法少女小圆》的一集，并以捷克剧作家、活动家和政治家瓦茨拉夫·哈维尔的引言片段作为标题，紧接着是该集的动画标题。我选择引用哈维尔原因也会随着本书的展开变得明了。书名同样取自哈维尔之言：“或许绝望正是滋养人类希望的土壤；或许人无法找到生活的意义，除非首先体验到其荒谬性。” \cite{ref9} 

第一部分之后是第一个插曲《完美一刻》，讨论该系列如何运用了歌德《浮士德》中的叙事元素。

第二部分是最短的部分，每章分别探讨在TV动画播出期间及之后发布的三部漫画衍生作品。这些章节的标题均取自TV动画中的一句台词。

接下来的第二个插曲《奶之尸》则探讨在剧场版《魔法少女小圆新篇：叛逆的物语》（以下简称《叛逆》）中，奶酪的含义与炼金术中腐化的概念。

最后一部分《叛逆》由七章组成，探讨《叛逆》这部电影。我只选择这部电影，是因为前两部电影基本上是对TV动画的总结，几乎没有新画面，更谈不上新剧情。与此不同的是，《叛逆》作为TV动画的续作\cite{ref10}。这七章每章讨论一个《叛逆》可能反叛的目标，因此每章标题的形式都是“叛逆的\underline{\makebox[1cm]{}}”。

本书是一部分析性的著作，而非导读。剧情、角色和制作细节只在分析相关内容时描述，一般情况下，本书假设读者对所讨论的作品已经熟悉。撰写本书时，该TV动画在美国可以通过DVD、Blu-ray光盘和Netflix与CrunchyRoll流媒体服务轻松观看；所讨论的漫画也已经出版，或正在美国以英文翻译出版。由于《叛逆》仅进行了短暂的影院上映，因此要看起来比较困难；不过可以购买附带英文字幕的日本版Blu-ray光盘，并且预计会在2015年4月发行带有英文字幕和配音选项的Blu-ray版。鉴于读者对以上这些作品熟悉的假设，本书不对所讨论的作品进行引用，仅在涉及《魔法少女小圆》系列之外的参考文献时进行引用。

最后，本书的一些章节，尤其是第2章和第9章，深入分析了动画中特别是整个魔法少女类型中的某些女性主义问题。在这一讨论中，请牢记所使用的女性主义理论来源于西方，基于欧洲的文化历史和假设。因此，虽然它适用于西方观众的观看体验，但不一定适用于日本文化——即该动画的创作文化和主要观众。因此，这些章节的讨论主要适用于从西方角度来看待该动画的视角，不一定能普遍适用于日本的文化体验。1

\paragraph{致谢}~{}

本书的完成离不开很多人的帮助。首先要感谢的是我的朋友兼封面设计师Viga Gadson，正是她最早向我介绍了《魔圆》。其次是Unnoun，他不断将我博客上关于该作品的文章发布到各种论坛，使这些文章成为我写作生涯中浏览量最高的内容之一，也促使我写了更多关于《魔圆》的文章，最终汇成了这本书。第三要感谢的是我的编辑Kit Paige，本书中学术部分的修饰都是她的努力，而任何遗留的错误都是因为我的固执，并非她的疏忽。第四要感谢MarkWatches.net的 Mark Oshiro和SFDebris.com的Chuck Sonnenburg，他们对动画每集的评论帮助我在自己的想法上获得了急需的批判性建议。最后但同样重要的是，感谢我的博客JedABlue.com上的读者们，他们对我最初稿件的深刻评论帮助我暴露了弱点，指出了需要澄清的部分，或提供了我未曾考虑的批评意见。特别要感谢Alan Jacobs Richardson（01d55），我们关于《叛逆》与E.T.A. Hoffmann《胡桃夹子和老鼠国王》 \cite{ref11} 的关系的长篇讨论启发了第20章和第21章中的许多内容。
